\chapter{Roteiro de três dias}

\section*{Dia 1 — Fès el-Bali}

\begin{itemize}
\item Bab Boujloud;
\item Talaa Kebira;
\item Madrasa Bou Inania;
\item Dar al-Magana;
\item Museu Nejjarine;
\item Fonte Nejjarine;
\item Al-Qarawiyyin;
\item Praça Seffarine;
\item souks e bairros de artesãos;
\item curtumes;
\item Jnan Sbil no final da tarde.
\end{itemize}

Este é o dia mais intenso de caminhada. O ideal é contratar um guia particular durante a manhã e deixar a tarde mais livre para voltar aos lugares que chamaram atenção.

\section*{Dia 2 — Fès judaica e Fès el-Jdid}

\begin{itemize}
\item Mellah;
\item Cemitério Judaico;
\item Sinagoga Ibn Danan;
\item Fès el-Jdid;
\item Palácio Real e Bab el-Makhzen;
\item Dar Batha, se estiver aberto;
\item caminhada pelo bairro;
\item Jnan Sbil, caso não tenha sido visitado no primeiro dia.
\end{itemize}

Este é o dia mais diretamente ligado à história judaica da cidade. Vale reservar tempo para caminhar pelo Mellah em vez de tratá-lo apenas como uma sequência de pontos turísticos.

\section*{Dia 3 — arquitetura, artesanato e panorama}

\begin{itemize}
\item Madrasa Al-Attarine;
\item Praça Seffarine;
\item oficinas de cerâmica;
\item madeira e outros artesanatos;
\item tempo livre na medina;
\item programação cultural ou exposição de manuscritos, se houver;
\item música andalusina à noite, se houver concerto;
\item Túmulos Merínidas no final da tarde.
\end{itemize}

O terceiro dia deve ser mais flexível. Se vocês já estiverem cansados da medina, vale reduzir o número de monumentos e dedicar mais tempo a cafés, jardins, oficinas e pequenos espaços culturais.

---

\chapter{Roteiro de dois dias}

\textbf{Dia 1 — Medina, porto e fortificações}

Começar pela Skala de la Ville e seguir pelas muralhas até a região da Kasbah. Caminhar pela medina, passando pelos souks e pelas oficinas de artesãos, até a Place Moulay Hassan. Depois, seguir para a Skala do Porto e Bab el-Marsa. Reservar algum tempo para o porto, especialmente no período da manhã. No fim da tarde, caminhar pela praia ou retornar às muralhas para observar o pôr do sol.

\textbf{Dia 2 — Mellah e memória judaica}

Começar pelo Mellah e percorrer suas ruas com atenção à arquitetura e à antiga organização do bairro. Visitar Bayt Dakira e, se estiver aberto, a Sinagoga Chaim Pinto. Seguir até o antigo cemitério judaico. Depois, visitar o Musée Sidi Mohammed Ben Abdellah e voltar à medina para explorar artesanato, galerias e pequenos mercados. À noite, procurar uma apresentação de música Gnawa ou outro programa cultural local.

%%%%%%%%%%%%%%%%%%%%%%%%%%%%%%%%%%%%%%%%%%%%%%%%%%%%%%%%%%%%%%%%%%%%%%%%%%%%%%%%%%%
